\section{冻结的 OLMo 中的单比特规则传递}
\label{sec:olmo}
在规则~110 实验中，构型合并后，观察者将原始传感器读数 $Z$ 写回一个元胞。自动机利用格点中保留的比特 $K$，在下一次更新时恢复过去的比特 $Y$。在 OLMo 的两步实验中，发送方在新比特串出现之前，将过去的选择 $h$ 和反馈 $e$ 编为一个规则词元 $m\in\{A,B\}$。接收方使用相同的冻结权重，从空的 KV 缓存开始新一轮计算，再将规则应用于输入串。以此划分两轮计算，可以测量已选定规则的传递。要考察模型在连续保留的内部状态中如何修正策略，还需另设实验。

\subsection{四种反事实条件}
设模型先前选择了 $h=0$（A，取首位）或 $h=1$（B，取末位）。选择正确时观察者返回 $e=0$，选错时返回 $e=1$。对每种选择都考察两个反馈分支。给出任何新串之前，修正后的规则为 $p=h\oplus e$。给出 $\mathbf{x}=LwR$ 后，正确答案依规则而定，为 $f_0(\mathbf{x})=L$ 或 $f_1(\mathbf{x})=R$。当首尾不同时，在相同的 $e,\mathbf{x}$ 下，具有不同 $h$ 的两段历史要求相反的动作；当首尾相同时，两者要求输出相同的比特。呈现首尾的全部四种组合，可以检验仅依靠反馈标记与答案之间的关联来作答的策略。

每个中间串对应一个区组，包含 $(L,R,h,e)$ 的 16 种条件。开发集（dev）包含测量前选定的八个中间串。只有全部 16/16 个答案正确，该区组才计为通过。通过门槛为 6/8 个区组。每次必需的数字读出均要求全词表上的概率质量满足 $P(0)+P(1)\ge0.80$。当 logit 差值 $\ell_1-\ell_0$ 非零时，其符号决定所选比特；零差值计作弃答。具有相同前缀的各行复用相应的模型计算。

\subsection{测量结果}
在 G0 中，模型处理完所选词元后，先隔离缓存，再进行诊断读出。在 G0R 中，用另一段自然产生的历史中对应位置的值替换选择位置的缓存，然后恢复原值。在最初的 Base 实验中，两个交换方向共 256/256 次完整 logit 向量比较均逐字节重现了原值，最大绝对差为零。这一对照验证回放的完整性。模型会否改选相反策略，还需行为检验。

\begin{center}\small
\begin{tabularx}{\linewidth}{@{}Xrrrl@{}}\toprule
检查点与接口（dev） & 条件行 & 区组 & $e=1$ 时的配对 & 状态 \\\midrule
Base，原始 raw 接口 & 76/128 & 0/8 & 0/16 & \texttt{FAIL\_DEV}\\
Instruct，对齐后的 raw 接口 & 79/128 & 0/8 & 1/16 & \texttt{FAIL\_FORMAT}\\
Instruct，chat，correct/wrong & 115/128 & 2/8 & 9/16 & \texttt{FAIL\_DEV}\\
Instruct，chat，keep/switch & 109/128 & 0/8 & 6/16 & \texttt{FAIL\_DEV}\\
Instruct，chat，两步（A） & 124/128 & 6/8 & 16/16 & \texttt{PASS\_DEV}\\
\bottomrule\end{tabularx}
\end{center}
表中一对通过，要求在相同 $e,\mathbf{x}$ 且 $L\ne R$ 时，对先前选择相反的两段历史均回答正确。前两行之间，权重和示例都发生了变化。在 Instruct 的 raw 接口中，显式给出规则的对照条件有 28/64 行未达到概率质量门槛。使用原生聊天模板和 correct/wrong 反馈时，在独立的缓存副本上正确读出了 2/2 个先前选择。显式给出规则时，答对了 62/64 行，通过了 6/8 个区组。Macar 等人向模型激活中注入向量，分别测量模型能否察觉干预、能否说出注入的概念~\cite{macar2026}。这里的行为判据要求模型把收到的信息应用于新串。激活干预实验 G3 的状态为 \texttt{NOT\_RUN}。

\subsection{单比特信道定理}
在两步方案中，发送方在 $\mathbf{x}$ 出现前输出全词表中得分唯一最高的词元。只有词元 A 和 B 是合格输出。接收方从自己的空 KV 缓存开始，只接收词元 $m$ 的 ID，重新计算其表示，随后才接收 $\mathbf{x}$。接收方无法读取发送方的分支标识、激活或 logits。外部程序不使用规则标签选择编码。

\begin{proposition}
设 $m$ 是接收方从上一轮获得的唯一可变输入，其字母表有两个元素，信道形如 $(h,e)\longrightarrow m\longrightarrow F(m,\mathbf{x})$。若在同一个满足 $L\ne R$ 的输入上，对四段历史都作出正确回答，则按 $00,01,10,11$ 的顺序，发送方的编码表只能是 $0110$ 或 $1001$。
\end{proposition}
\begin{proof}
在给定 $\mathbf{x}$ 时，$p$ 的两个不同取值要求不同的答案。若两段具有不同 $p$ 的历史对应同一 $m$，接收方就会在目标不同的情况下得到相同的输入。因此两组非空的消息集合必不相交。总共只有两个消息，故每个 $p$ 恰好对应一个 $m$。两种可能只相差词元名称的置换。
\end{proof}

在 dev 上观察到的表为 $A,B,B,A$。对每一个 $e$，在首尾不同的 16/16 对条件中，两段历史均答对。传递消息时得到 64/64 个正确答案，独立计算的恒定 A 对照得到 32/64。发送方的四次选择在接收方处理各个新串时复用。由此可知，消息在输入串出现前已经携带规则类别的信息。要确定发送方内部如何计算该信息，还需要其他干预。多维状态则可能同时传递 $(h,e)$ 两个比特，将规则的计算留给接收方。测量来源和对照见附录~\ref{app:gates}。
